\IfFileExists{stacks-project.cls}{%
\documentclass{stacks-project}
}{%
\documentclass{amsart}
}

% For dealing with references we use the comment environment
\usepackage{verbatim}
\newenvironment{reference}{\comment}{\endcomment}
%\newenvironment{reference}{}{}
\newenvironment{slogan}{\comment}{\endcomment}
\newenvironment{history}{\comment}{\endcomment}

% For commutative diagrams we use Xy-pic
\usepackage[all]{xy}

% We use 2cell for 2-commutative diagrams.
\xyoption{2cell}
\UseAllTwocells

% We use multicol for the list of chapters between chapters
\usepackage{multicol}

% This is generally recommended for better output
\usepackage{lmodern}
\usepackage[T1]{fontenc}

% For cross-file-references
\usepackage{xr-hyper}

% Package for hypertext links:
\usepackage{hyperref}

% For any local file, say "hello.tex" you want to link to please
% use \externaldocument[hello-]{hello}
\externaldocument[introduction-]{introduction}
\externaldocument[conventions-]{conventions}
\externaldocument[sets-]{sets}
\externaldocument[categories-]{categories}
\externaldocument[topology-]{topology}
\externaldocument[sheaves-]{sheaves}
\externaldocument[sites-]{sites}
\externaldocument[stacks-]{stacks}
\externaldocument[fields-]{fields}
\externaldocument[algebra-]{algebra}
\externaldocument[brauer-]{brauer}
\externaldocument[homology-]{homology}
\externaldocument[derived-]{derived}
\externaldocument[simplicial-]{simplicial}
\externaldocument[more-algebra-]{more-algebra}
\externaldocument[smoothing-]{smoothing}
\externaldocument[modules-]{modules}
\externaldocument[sites-modules-]{sites-modules}
\externaldocument[injectives-]{injectives}
\externaldocument[cohomology-]{cohomology}
\externaldocument[sites-cohomology-]{sites-cohomology}
\externaldocument[dga-]{dga}
\externaldocument[dpa-]{dpa}
\externaldocument[sdga-]{sdga}
\externaldocument[hypercovering-]{hypercovering}
\externaldocument[schemes-]{schemes}
\externaldocument[constructions-]{constructions}
\externaldocument[properties-]{properties}
\externaldocument[morphisms-]{morphisms}
\externaldocument[coherent-]{coherent}
\externaldocument[divisors-]{divisors}
\externaldocument[limits-]{limits}
\externaldocument[varieties-]{varieties}
\externaldocument[topologies-]{topologies}
\externaldocument[descent-]{descent}
\externaldocument[perfect-]{perfect}
\externaldocument[more-morphisms-]{more-morphisms}
\externaldocument[flat-]{flat}
\externaldocument[groupoids-]{groupoids}
\externaldocument[more-groupoids-]{more-groupoids}
\externaldocument[etale-]{etale}
\externaldocument[chow-]{chow}
\externaldocument[intersection-]{intersection}
\externaldocument[pic-]{pic}
\externaldocument[weil-]{weil}
\externaldocument[adequate-]{adequate}
\externaldocument[dualizing-]{dualizing}
\externaldocument[duality-]{duality}
\externaldocument[discriminant-]{discriminant}
\externaldocument[derham-]{derham}
\externaldocument[local-cohomology-]{local-cohomology}
\externaldocument[algebraization-]{algebraization}
\externaldocument[curves-]{curves}
\externaldocument[resolve-]{resolve}
\externaldocument[models-]{models}
\externaldocument[functors-]{functors}
\externaldocument[equiv-]{equiv}
\externaldocument[pione-]{pione}
\externaldocument[etale-cohomology-]{etale-cohomology}
\externaldocument[proetale-]{proetale}
\externaldocument[relative-cycles-]{relative-cycles}
\externaldocument[more-etale-]{more-etale}
\externaldocument[trace-]{trace}
\externaldocument[crystalline-]{crystalline}
\externaldocument[spaces-]{spaces}
\externaldocument[spaces-properties-]{spaces-properties}
\externaldocument[spaces-morphisms-]{spaces-morphisms}
\externaldocument[decent-spaces-]{decent-spaces}
\externaldocument[spaces-cohomology-]{spaces-cohomology}
\externaldocument[spaces-limits-]{spaces-limits}
\externaldocument[spaces-divisors-]{spaces-divisors}
\externaldocument[spaces-over-fields-]{spaces-over-fields}
\externaldocument[spaces-topologies-]{spaces-topologies}
\externaldocument[spaces-descent-]{spaces-descent}
\externaldocument[spaces-perfect-]{spaces-perfect}
\externaldocument[spaces-more-morphisms-]{spaces-more-morphisms}
\externaldocument[spaces-flat-]{spaces-flat}
\externaldocument[spaces-groupoids-]{spaces-groupoids}
\externaldocument[spaces-more-groupoids-]{spaces-more-groupoids}
\externaldocument[bootstrap-]{bootstrap}
\externaldocument[spaces-pushouts-]{spaces-pushouts}
\externaldocument[spaces-chow-]{spaces-chow}
\externaldocument[groupoids-quotients-]{groupoids-quotients}
\externaldocument[spaces-more-cohomology-]{spaces-more-cohomology}
\externaldocument[spaces-simplicial-]{spaces-simplicial}
\externaldocument[spaces-duality-]{spaces-duality}
\externaldocument[formal-spaces-]{formal-spaces}
\externaldocument[restricted-]{restricted}
\externaldocument[spaces-resolve-]{spaces-resolve}
\externaldocument[formal-defos-]{formal-defos}
\externaldocument[defos-]{defos}
\externaldocument[cotangent-]{cotangent}
\externaldocument[examples-defos-]{examples-defos}
\externaldocument[algebraic-]{algebraic}
\externaldocument[examples-stacks-]{examples-stacks}
\externaldocument[stacks-sheaves-]{stacks-sheaves}
\externaldocument[criteria-]{criteria}
\externaldocument[artin-]{artin}
\externaldocument[quot-]{quot}
\externaldocument[stacks-properties-]{stacks-properties}
\externaldocument[stacks-morphisms-]{stacks-morphisms}
\externaldocument[stacks-limits-]{stacks-limits}
\externaldocument[stacks-cohomology-]{stacks-cohomology}
\externaldocument[stacks-perfect-]{stacks-perfect}
\externaldocument[stacks-introduction-]{stacks-introduction}
\externaldocument[stacks-more-morphisms-]{stacks-more-morphisms}
\externaldocument[stacks-geometry-]{stacks-geometry}
\externaldocument[moduli-]{moduli}
\externaldocument[moduli-curves-]{moduli-curves}
\externaldocument[examples-]{examples}
\externaldocument[exercises-]{exercises}
\externaldocument[guide-]{guide}
\externaldocument[desirables-]{desirables}
\externaldocument[coding-]{coding}
\externaldocument[obsolete-]{obsolete}
\externaldocument[fdl-]{fdl}
\externaldocument[index-]{index}

% Theorem environments.
%
\theoremstyle{plain}
\newtheorem{theorem}[subsection]{Theorem}
\newtheorem{proposition}[subsection]{Proposition}
\newtheorem{lemma}[subsection]{Lemma}

\theoremstyle{definition}
\newtheorem{definition}[subsection]{Definition}
\newtheorem{example}[subsection]{Example}
\newtheorem{exercise}[subsection]{Exercise}
\newtheorem{situation}[subsection]{Situation}

\theoremstyle{remark}
\newtheorem{remark}[subsection]{Remark}
\newtheorem{remarks}[subsection]{Remarks}

\numberwithin{equation}{subsection}

% Macros
%
\def\lim{\mathop{\mathrm{lim}}\nolimits}
\def\colim{\mathop{\mathrm{colim}}\nolimits}
\def\Spec{\mathop{\mathrm{Spec}}}
\def\Hom{\mathop{\mathrm{Hom}}\nolimits}
\def\Ext{\mathop{\mathrm{Ext}}\nolimits}
\def\SheafHom{\mathop{\mathcal{H}\!\mathit{om}}\nolimits}
\def\SheafExt{\mathop{\mathcal{E}\!\mathit{xt}}\nolimits}
\def\Sch{\mathit{Sch}}
\def\Mor{\mathop{\mathrm{Mor}}\nolimits}
\def\Ob{\mathop{\mathrm{Ob}}\nolimits}
\def\Sh{\mathop{\mathit{Sh}}\nolimits}
\def\NL{\mathop{N\!L}\nolimits}
\def\CH{\mathop{\mathrm{CH}}\nolimits}
\def\proetale{{pro\text{-}\acute{e}tale}}
\def\etale{{\acute{e}tale}}
\def\QCoh{\mathit{QCoh}}
\def\Ker{\mathop{\mathrm{Ker}}}
\def\Im{\mathop{\mathrm{Im}}}
\def\Coker{\mathop{\mathrm{Coker}}}
\def\Coim{\mathop{\mathrm{Coim}}}

% Boxtimes
%
\DeclareMathSymbol{\boxtimes}{\mathbin}{AMSa}{"02}

%
% Macros for moduli stacks/spaces
%
\def\QCohstack{\mathcal{QC}\!\mathit{oh}}
\def\Cohstack{\mathcal{C}\!\mathit{oh}}
\def\Spacesstack{\mathcal{S}\!\mathit{paces}}
\def\Quotfunctor{\mathrm{Quot}}
\def\Hilbfunctor{\mathrm{Hilb}}
\def\Curvesstack{\mathcal{C}\!\mathit{urves}}
\def\Polarizedstack{\mathcal{P}\!\mathit{olarized}}
\def\Complexesstack{\mathcal{C}\!\mathit{omplexes}}
% \Pic is the operator that assigns to X its picard group, usage \Pic(X)
% \Picardstack_{X/B} denotes the Picard stack of X over B
% \Picardfunctor_{X/B} denotes the Picard functor of X over B
\def\Pic{\mathop{\mathrm{Pic}}\nolimits}
\def\Picardstack{\mathcal{P}\!\mathit{ic}}
\def\Picardfunctor{\mathrm{Pic}}
\def\Deformationcategory{\mathcal{D}\!\mathit{ef}}
\usepackage{cleveref}
\crefname{originallemma}{original lemma}{original lemmas}
\Crefname{originallemma}{Original lemma}{Original lemmas}
\makeatletter
\newlabel{lemma-adequate-flat}{{06UW}{source}{}{sourcepointer-1}{}}
\newlabel{lemma-adequate-flat@cref}{{[originallemma][0][]06UW}{[1][0][]source}{}{}{}}
\makeatother
\crefname{originallemma}{original lemma}{original lemmas}
\Crefname{originallemma}{Original lemma}{Original lemmas}
\makeatletter
\newlabel{lemma-flat-functor-split}{{06UZ}{source}{}{sourcepointer-2}{}}
\newlabel{lemma-flat-functor-split@cref}{{[originallemma][0][]06UZ}{[1][0][]source}{}{}{}}
\makeatother
\crefname{originallemma}{original lemma}{original lemmas}
\Crefname{originallemma}{Original lemma}{Original lemmas}
\makeatletter
\newlabel{lemma-flat-linear-functor}{{06V5}{source}{}{sourcepointer-3}{}}
\newlabel{lemma-flat-linear-functor@cref}{{[originallemma][0][]06V5}{[1][0][]source}{}{}{}}
\makeatother
\title{An extension of a linearly adequate functor by a module-valued functor is adequate}
\author{The Stacks Project Authors}
\date{stacks-project:adequate.tex:local-snapshot-20261012; GNU FDL 1.2 or later}
\begin{document}
\maketitle
\noindent\textit{Editorial note.} All delivered mathematical text is the author's own, copied from the pinned Stacks Project chapter file adequate.tex; no statement or proof step was rewritten, reordered or removed.\par\medskip
\noindent\textit{Editorial note.} The source defines the three algebra maps p\_1, p\_2, m by p\_1(x\_i) = x\_i, p\_1(x\_i) = y\_i, m(x\_i) = x\_i + y\_i -- that is, the line for p\_2 is printed as p\_1 in the frozen source. The text is preserved verbatim and the observation is recorded here rather than corrected.\par\medskip
\noindent\textit{Editorial note.} The counted spans carry navigation stubs for the references leaving them; no bibliography citation occurs in the spans.\par\medskip
\noindent\textit{Editorial note.} The preamble span points at the chapter's own preamble.tex shipped with the frozen source, used verbatim as staged.\par\medskip
\noindent\textit{Allowed standard prerequisites (editorial).}\begin{itemize}
\item The flat-base-change property for M-underlined and for linearly adequate functors. Same chapter, Lemma lemma-adequate-flat, invoked at proof lines 791-793. Together with the five lemma it gives G(B) tensor\_B B' = G(B') for flat B -> B'.
\item A module-valued functor with the flat-base-change property splits into weight spaces over a graded algebra. Same chapter, Lemma lemma-flat-functor-split, invoked at proof lines 793-796. Gives G(B) = direct sum of its graded pieces, allowing the lift of xi.
\item The structure of linearly adequate functors and their presentation. Same chapter, Lemma lemma-flat-linear-functor, invoked within the proof's use of the presentation of L. Lets the author describe L(C) as the solutions of the relations and build the graded algebras B and C.
\end{itemize}\medskip
\ifdefined\hypertarget\hypertarget{sourcepointer-1}{}\fi
\noindent\textit{Source reference.} Original lemma 06UW: sources/stacks-adequate/pinned-original-source.tex, source lines 442--442. Pointer outside the counted span, invoked once for context; the load-bearing argument of the counted statement is the author's proof reproduced in full in the proof span.\par\medskip
\ifdefined\hypertarget\hypertarget{sourcepointer-2}{}\fi
\noindent\textit{Source reference.} Original lemma 06UZ: sources/stacks-adequate/pinned-original-source.tex, source lines 533--533. Pointer outside the counted span, invoked once for context; the load-bearing argument of the counted statement is the author's proof reproduced in full in the proof span.\par\medskip
\ifdefined\hypertarget\hypertarget{sourcepointer-3}{}\fi
\noindent\textit{Source reference.} Original lemma 06V5: sources/stacks-adequate/pinned-original-source.tex, source lines 750--750. Pointer outside the counted span, invoked once for context; the load-bearing argument of the counted statement is the author's proof reproduced in full in the proof span.\par\medskip
\begin{lemma}
\label{lemma-extension-adequate-key}
Let $A$ be a ring.
Let $0 \to \underline{M} \to G \to L \to 0$ be a short exact sequence
of module-valued functors on $\textit{Alg}_A$ with $L$ linearly adequate.
Then $G$ is adequate.
\end{lemma}

\begin{proof}
We first point out that for any flat $A$-algebra map
$B \to B'$ the map $G(B) \otimes_B B' \to G(B')$ is an isomorphism.
Namely, this holds for $\underline{M}$ and $L$, see
Lemma \ref{lemma-adequate-flat}
and hence follows for $G$ by the five lemma. In particular, by
Lemma \ref{lemma-flat-functor-split}
we see that $G(B) = \bigoplus_{k \in \mathbf{Z}} G(B)^{(k)}$
for any graded $A$-algebra $B$.

\medskip\noindent
Choose an exact sequence
$0 \to L \to \underline{A^{\oplus n}} \to \underline{A^{\oplus m}}$.
Suppose that the map $A^{\oplus n} \to A^{\oplus m}$ is given by the
$m \times n$-matrix $(a_{ij})$. Consider the graded $A$-algebra
$B = A[x_1, \ldots, x_n]/(\sum a_{ij}x_j)$. The element
$(x_1, \ldots, x_n) \in \underline{A^{\oplus n}}(B)$ maps to zero in
$\underline{A^{\oplus m}}(B)$ hence is the image of a unique element
$\xi \in L(B)$. Observe that $\xi \in L(B)^{(1)}$. The map
$$
\Hom_A(B, C) \longrightarrow L(C), \quad
f \longmapsto L(f)(\xi)
$$
defines an isomorphism of functors. The reason is that $f$ is
determined by the images $c_i = f(x_i) \in C$ which have to
satisfy the relations $\sum a_{ij}c_j = 0$. And $L(C)$ is the
set of $n$-tuples $(c_1, \ldots, c_n)$ satisfying the relations
$\sum a_{ij} c_j = 0$.

\medskip\noindent
Since the value of each of the functors $\underline{M}$, $G$, $L$
on $B$ is a direct sum of its weight spaces (by the lemma mentioned
above) exactness of $0 \to \underline{M} \to G \to L \to 0$ implies
the sequence $0 \to \underline{M}(B)^{(1)} \to G(B)^{(1)} \to L(B)^{(1)} \to 0$
is exact. Thus we may choose an element $\theta \in G(B)^{(1)}$ mapping
to $\xi$.

\medskip\noindent
Consider the graded $A$-algebra
$$
C = A[x_1, \ldots, x_n, y_1, \ldots, y_n]/
(\sum a_{ij}x_j, \sum a_{ij}y_j)
$$
There are three graded $A$-algebra homomorphisms $p_1, p_2, m : B \to C$
defined by the rules
$$
p_1(x_i) = x_i, \quad
p_1(x_i) = y_i, \quad
m(x_i) = x_i + y_i.
$$
We will show that the element
$$
\tau = G(m)(\theta) - G(p_1)(\theta) - G(p_2)(\theta) \in G(C)
$$
is zero. First, $\tau$ maps to zero in $L(C)$ by a direct calculation.
Hence $\tau$ is an element of $\underline{M}(C)$.
Moreover, since $m$, $p_1$, $p_2$ are graded algebra maps we see
that $\tau \in G(C)^{(1)}$ and since $\underline{M} \subset G$
we conclude
$$
\tau \in \underline{M}(C)^{(1)} = M \otimes_A C_1.
$$
We may write uniquely
$\tau = \underline{M}(p_1)(\tau_1) + \underline{M}(p_2)(\tau_2)$ with
$\tau_i \in M \otimes_A B_1 = \underline{M}(B)^{(1)}$ because
$C_1 = p_1(B_1) \oplus p_2(B_1)$.
Consider the ring map $q_1 : C \to B$ defined by $x_i \mapsto x_i$ and
$y_i \mapsto 0$. Then
$\underline{M}(q_1)(\tau) =
\underline{M}(q_1)(\underline{M}(p_1)(\tau_1) + \underline{M}(p_2)(\tau_2)) =
\tau_1$.
On the other hand, because
$q_1 \circ m = q_1 \circ p_1$ we see that
$G(q_1)(\tau) = - G(q_1 \circ p_2)(\tau)$. Since $q_1 \circ p_2$ factors as
$B \to A \to B$ we see that $G(q_1 \circ p_2)(\tau)$ is in
$G(B)^{(0)}$, see
Lemma \ref{lemma-flat-functor-split}.
Hence $\tau_1 = 0$ because it is in
$G(B)^{(0)} \cap \underline{M}(B)^{(1)} \subset
G(B)^{(0)} \cap G(B)^{(1)} = 0$.
Similarly $\tau_2 = 0$, whence $\tau = 0$.

\medskip\noindent
Since $\theta \in G(B)$ we obtain a transformation of functors
$$
\psi : L(-) = \Hom_A(B, - ) \longrightarrow G(-)
$$
by mapping $f : B \to C$ to $G(f)(\theta)$. Since $\theta$ is a lift of
$\xi$ the map $\psi$ is a right inverse of $G \to L$. In terms of
$\psi$ the statements proved above have the following meaning:
$\tau = 0$ means that $\psi$ is additive and
$\theta \in G(B)^{(1)}$ implies that for any $A$-algebra $D$ we have
$\psi(ul) = u\psi(l)$ in $G(D)$ for $l \in L(D)$ and $u \in D^*$ a unit.
This implies that $\psi$ is a morphism of module-valued functors, see
Lemma \ref{lemma-flat-linear-functor}.
Clearly this implies that $G \cong \underline{M} \oplus L$ and we win.
\end{proof}

\end{document}
